\section{Introduction}

\subsection{Background and main result}

The study of descent sets for permutations may be traced back to Euler. 
A cyclic extension of this classical concept was introduced in the study of
Lie algebras~\cite{Klyachko} and descent algebras~\cite{Cellini}. Surprising connections of the cyclic descent
notion to a variety of mathematical areas were found later.
	
The {\em descent set} of a permutation $\pi = [\pi_1, \ldots, \pi_n]$ in the symmetric group $\symm_n$ on $n$ letters is
\[
    \Des(\pi) := \{1 \le i \le n-1 \,:\, \pi_i > \pi_{i+1} \}
    \quad \subseteq [n-1],
\]
where $[m]:=\{1,2,\ldots,m\}$.  
Cellini~\cite{Cellini} introduced a natural notion of {\em cyclic descent set}: 
\[
    \CDes(\pi) := \{1 \leq i \leq n \,:\, \pi_i > \pi_{i+1} \}
    \quad \subseteq [n],
\]
with the convention $\pi_{n+1}:=\pi_1$.  
The more restricted notion of cyclic descent number had been used previously by Klyachko~\cite{Klyachko}.
This cyclic descent set was further studied by Dilks, Petersen and Stembridge~\cite{DPS} and others.
	
There exists a well-established notion of descent set for standard Young tableaux (SYT), but it has no obvious cyclic analogue. 
In a breakthrough work, Rhoades~\cite{Rhoades} defined a notion of cyclic descent set for standard Young tableaux of rectangular shape.
The properties common to Cellini's definition (for permutations) and Rhoades' construction (for SYT) appeared in other combinatorial settings as well~\cite{PPR, Pechenik, ER, AER}.
This led to an abstract definition, as follows.
	
\begin{definition}[\cite{ARR}]\label{def:cDes}
Let $\TTT$ be a finite set, equipped with a set valued map (called {\em descent map}) 
$\Des: \TTT \longrightarrow 2^{[n-1]}$. 
Let ${\rm shift}: 2^{[n]} \longrightarrow 2^{[n]}$ be the mapping on subsets of~$[n]$ induced by the cyclic shift $i \mapsto i + 1 \pmod n$ of elements $i \in [n]$.
A {\em cyclic extension} of $\Des$ is
a pair $(\cDes,p)$, where 
$\cDes: \TTT \longrightarrow 2^{[n]}$ is a map 
and $p: \TTT \longrightarrow \TTT$ is a bijection,
satisfying the following axioms:  for all $T$ in  $\TTT$,
\[
\begin{array}{rl}
\text{(extension)}   & \cDes(T) \cap [n-1] = \Des(T),\\
\text{(equivariance)}& \cDes(p(T))  = {\rm shift}(\cDes(T)),\\
\text{(non-Escher)}  & \varnothing \subsetneq \cDes(T) \subsetneq [n].\\
\end{array}
\]
\end{definition}

The term ``non-Escher'' refers to M.\ C.\ Escher's drawing ``Ascending and Descending'', which illustrates the impossibility
of the cases $\cDes(\pi) = \varnothing$ and $\cDes(\pi) = [n]$ for permutations $\pi \in S_n$.

A {\em ribbon} is a skew shape which contains no $2 \times 2$ square.
	
\begin{theorem}[{\cite[Theorem 1.1]{ARR}}]\label{conj1}
    Let $\lambda/\mu$ be a skew shape with $n$ cells.
    The descent map $\Des$ on $\SYT(\lambda/\mu)$ has a cyclic extension $(\cDes,p)$ if and only if~$\lambda/\mu$ is not a connected ribbon.
\end{theorem}
	
The original proof used Postnikov's toric symmetric functions.
A constructive proof was recently given by Huang~\cite{Huang}.

For connections of cyclic descents to 
Kazhdan--Lusztig theory see~\cite{Rhoades}; for topological aspects and connections to the  
Steinberg torus see~\cite{DPS}; for twisted Sch\"utzenberger promotion see~\cite{Rhoades, Huang}; for cyclic quasisymmetric functions and Schur-positivity see~\cite{AGRR};	for Postnikov's toric Schur functions see~\cite{ARR}.
The goal of this paper is to determine which conjugacy classes of the symmetric group carry a cyclic descent extension.
		
\begin{observation}
    Let $\Des$ and $\CDes$ denote the classical descent set and Cellini's cyclic descent set on permutations, respectively.
    Let $p: S_n \rightarrow S_n$ be the rotation 
    \[
        \begin{array}{rcl}
            [\pi_1,\pi_2,\ldots,\pi_{n-1},\pi_n]&\overset{p}{\longmapsto}&
            [\pi_n,\pi_1,\pi_2,\ldots,\pi_{n-1}]. 
        \end{array}
    \]
    Then the pair $(\CDes,p)$ is a cyclic descent extension of $\Des$ on $S_n$ in the sense of Definition~\ref{def:cDes}.
\end{observation}
	
Cellini's definition provides a cyclic extension of $\Des$ on the full symmetric group, but not on some of its subsets -- for example, on many conjugacy classes;
see Section~\ref{sec:Cellini} below.

\begin{example}
    Consider the conjugacy class of 4-cycles in $S_4$,
    \[
        \C_{(4)}=\{2341, 2413, 3142, 3421, 4123, 4312\}.
    \]
    Cellini's cyclic descent sets are
    \[
        \{3\}, \{2,4\}, \{1,3\}, \{2,3\}, \{1\}, \{1,2\},
    \]
    respectively; this family is not closed under cyclic rotation.
    On the other hand, redefining the cyclic descent sets to be
    \[
        \cDes(2341)=\{3,4\},\ \cDes(2413)=\{2,4\}, \ \cDes(3142)=\{1,3\}, 
    \]
    \[
        \cDes(3421)=\{2,3\},\ \cDes(4123)=\{1,4\},\ \cDes(4312)=\{1,2\}
    \]
    and defining the map $p$ by
    \[
        2341 \rightarrow 4123 \rightarrow 4312 \rightarrow 3421 \rightarrow 2341
    \]
    and
    \[
        3142 \rightarrow 2413 \rightarrow 3142,
    \]
    the pair $(\cDes,p)$ does determine a cyclic extension of $\Des$ for this conjugacy class.
\end{example}
	
The goal of this paper is to show that most conjugacy classes in $S_n$ carry a cyclic descent extension.
In fact, we obtain a full characterization.

Recall that an integer is \emph{square-free} if no prime square divides it; in particular, $1$ is square-free. Our main result is the following. 
	
\begin{theorem}\label{thm:main}
    Let $\lambda$ be a partition of $n$, and let
    $\C_\lambda \subseteq S_n$ be the corresponding conjugacy class. 
    The descent map $\Des$ on $\C_\lambda$ has a cyclic extension $(\cDes,p)$ if and only if~$\lambda$ is not of the form $(r^s)$ for some square-free $r$.
\end{theorem}
	

\subsection{Proof method} 
	
The proof of Theorem~\ref{thm:main} is non-constructive and   
involves a detailed study of the hook constituents in higher Lie characters.
Here is an overview of the main ingredients.


\subsubsection{Higher Lie characters}

\begin{definition}\label{def:higher}
For a partition $\lambda$ of $n$, let $\C_\lambda$ be the conjugacy class consisting of all the permutations in $S_n$ of cycle type $\lambda$,
and let $\chi^\lambda$ denote the irreducible $S_n$-character corresponding to $\lambda$.
Let $Z_\lambda$ be the centralizer of a permutation in $\C_\lambda$ (defined up to conjugacy). 
If $k_i$ denotes the number of parts of $\lambda$ equal to $i$, then $Z_\lambda$ is isomorphic to the direct	product $\times_{i=1}^n \ZZ_i\wr S_{k_i}$.
Here and in the rest of the paper $\ZZ_i$ denotes the cyclic group of order $i$, using additive notation (integers modulo $i$).
	
For each $i$, let $\omega_i$ be the linear character on $\ZZ_i\wr S_{k_i}$ indexed by the
$i$-tuple of partitions $(\emptyset, (k_i), \emptyset,
\dots,\emptyset)$. 
In other words, let $\zeta_i$ be a primitive linear character on the cyclic group $\ZZ_i$, and extend it to the wreath product $\ZZ_i\wr S_{k_i}$ so that it is $\zeta_i^{k_i}$
on the base subgroup $\ZZ_i^{k_i}$ and trivial on the wreathing subgroup $S_{k_i}$. Denote this extension by~$\omega_i$.
Now let
\[
\omega^\lambda := \bigotimes_{i=1}^n \omega_i,
\]
a linear character on $Z_\lambda$. 
Define the corresponding {\em higher Lie character} to be the induced character
\[
\psi^\lambda := \omega^\lambda\uparrow_{Z_\lambda}^{S_n}.
\]
\end{definition}

The study of higher Lie characters can be traced back 
to Schur~\cite{schur}. 
An old problem of Thrall~\cite{thrall} 
is to provide an explicit combinatorial interpretation of the 
multiplicities of the irreducible characters in the higher Lie character,
see also~\cite[Exercise~7.89(i)]{EC2}.
Only partial results are known: the case $\lambda = (n)$ was solved by Kra\'skiewicz and Weyman~\cite{KraskiewiczWeyman};
D\'esarm\'enien and Wachs~\cite{DesarmenienWachs} resolved a coarser version of Thrall's problem 
for the sum of higher Lie characters over all  
derangements, see also~\cite{ReinerWebb}. 
The best result so far is
Schocker's expansion~\cite[Theorem~3.1]{schocker}, which however involves signs and rational coefficients. 
For recent discussions see, e.g., \cite{Reiner, AS, Sundaram}.

A remarkable theorem of  Gessel and Reutenauer~\cite[Theorem 2.1]{GR} 
applies higher Lie characters to describe the fiber sizes of the descent set map on conjugacy classes.  
Their proof applies an interpretation of higher Lie character $\psi^\lambda$ in terms of quasisymmetric functions 
(Theorem~\ref{thm:GR22} below). 
It follows that 
higher Lie characters can be used to prove the existence of cyclic descent extensions 
as explained below. 


\subsubsection{Hook multiplicities and cyclic descent extensions} 

Recall the standard notation~$s_\lambda$ for the Schur function indexed by a partition $\lambda$, 
as well as $\f_{n,D}$ for the fundamental quasisymmetric function indexed by a subset $D\subseteq [n-1]$;
see Definition~\ref{def:fundamental_QSF}.
A subset $\A\subseteq S_n$ is {\em Schur-positive} 
if the
associated quasisymmetric function
\[
\Q(\A):=\sum\limits_{a\in \A}\f_{n,\Des(a)}
\]
is symmetric and Schur-positive.

For an integer $0 \le k<n$ and a Schur-positive subset $\A \subseteq S_n$ denote
\[
m_{k,\A} := \langle \Q(\A), s_{(n-k,1^k)} \rangle,
\]
where $s_{(n-k,1^k)}$ is the Schur function indexed by the hook partition $(n-k,1^k)$.

First we prove the following key lemma, which provides an algebraic criterion for the existence of a cyclic descent extension. 

\begin{lemma}\label{lem:21}
A Schur-positive set $\A\subseteq S_n$ has a cyclic descent extension if and only if 
the following two conditions hold:
\[
\begin{array}{rl}
	\text{\rm (divisibility)}   & 
	\text{the polynomial }
	\sum_{k=0}^{n-1} m_{k,\A} x^k
	\text{ is divisible by } 1+x; \\
	\text{\rm (non-negativity)} & 
	\text{the quotient has nonnegative coefficients.}
\end{array}
\]
\end{lemma}		

See Lemma~\ref{lem:2}  below.


\subsubsection{Divisibility} 

By the 
Gessel--Reutenauer theorem, 
for every conjugacy class $\C_\lambda$ 
the quasisymmetric function 
$\Q(\C_\lambda)$ is the Frobenius image of the higher Lie character~$\psi^\lambda$, thus  $\C_\lambda$ is  Schur-positive; 
see Theorem~\ref{thm:GR22} below.

For a partition $\lambda \vdash n$ denote
\begin{equation}\label{notation:m_k_hLc}
m_{k,\lambda} := m_{k,\C_\lambda}
= \langle \Q(\C_\lambda), s_{(n-k,1^k)} \rangle
= \langle \psi^\lambda, \chi^{(n-k,1^k)} \rangle.
\end{equation}
    
\begin{proposition}\label{prop:divisibility} 
    The hook-multiplicity generating function of the higher Lie character $\psi^\lambda$
    \[
        M_\lambda(x):=\sum\limits_{k=0}^{n-1}m_{k,\lambda}x^k
    \]
    is divisible by $1+x$ 
    if and only if $\lambda$ is not of the form $(r^s)$ for a square-free integer $r$.
\end{proposition}

This divisibility condition is proved using an explicit evaluation of the higher Lie character on $n$-cycles; see Section~\ref{sec:divisibility} below.

    
\subsubsection{Non-negativity} 

In order to prove Theorem~\ref{thm:main}, it remains to show that the coefficients of the quotient  $M_\lambda(x)/(1+x)$ are nonnegative, whenever $\lambda$ is not of the form $(r^s)$ for a square-free $r$. 
It turns out that partitions $\lambda$ which have at least two different parts, namely not of the form $(r^s)$ for any $r$, are the easiest to handle. 
In that case, a factorization of the associated higher Lie character $\psi^\lambda$ is applied to prove the following. 

\begin{lemma}\label{lem:distinct}
    Let $\lambda = \mu \sqcup \nu$ be a disjoint union of nonempty partitions with no common part. 
    Then
    \begin{equation}
        \frac{M_\lambda(x)}{1+x} = M_\mu(x) M_\nu(x),     
    \end{equation}
    and its coefficients are thus non-negative.
\end{lemma}	

The core of the proof of Theorem~\ref{thm:main} is 
the case of $\lambda = (r^s)$.
For a fixed positive integer $r$, consider the formal power series
\[
    M_r(x,y) := \sum_{\substack{i \ge 0 \\ s \ge 1}} m_{i,(r^s)} x^i y^s,
\]
where $m_{i,(r^s)}$ is the multiplicity of the hook character $\chi^{(rs-i,1^i)}$ in the higher Lie character $\psi^{(r^s)}$. 
The following theorem completes the proof of Theorem~\ref{thm:main}.
 
\begin{theorem}\label{t:main_hook-alternating2}
    If $r$ is not square-free then the formal power series 
    \[
        \frac{M_r(x,y)}{1+x}
    \] 
	has non-negative integer coefficients.
\end{theorem}

\begin{proof}[Proof of Theorem~\ref{thm:main}.]
    Combine Lemma~\ref{lem:21}, Proposition~\ref{prop:divisibility}, Lemma~\ref{lem:distinct} and Theorem~\ref{t:main_hook-alternating2}. 
\end{proof}
    

\subsection{Outline}

The rest of the paper is organized as follows. 

Necessary background, including a cyclic analogue of the Gessel--Reutenauer Theorem, is given in Section~\ref{sec:background}.

The role of hooks in the study of cyclic descent extensions is explained in Section~\ref{sec:hooks}.
In particular, a necessary and sufficient criterion
for Schur-positive sets to carry a cyclic descent extension (Lemma~\ref{lem:21}) is proved in Subsection~\ref{sec:hook_criterion}.
Using this criterion, the proof of Theorem~\ref{thm:main} is reduced
to Proposition~\ref{prop:divisibility} (divisibility) and  Theorem~\ref{t:main_hook-alternating2} (non-negativity).
Proposition~\ref{prop:divisibility} is proved in Subsection~\ref{sec:divisibility}.  

The proof of Theorem~\ref{t:main_hook-alternating2}, stating the non-negativity of the coefficients of~$M_\lambda(x)/(1+x)$ whenever this quotient is a polynomial, spans 
Sections~\ref{sec:distinct}, \ref{sec:cycles} and~\ref{sec:Witt}:
the case of more than one cycle length is considered in Section~\ref{sec:distinct};
the case of a single cycle 
is considered in Section~\ref{sec:cycles};
and the case of cycle type $(r^s)$ with $s>1$ is considered in Section~\ref{sec:Witt}.
In the case of more than one cycle length, 
non-negativity is proved using a factorization of the associated higher Lie character (Lemma~\ref{lem:distinct}).  
In the case of a single cycle,  
combining a combinatorial formula for inner products with a variant of the Witt transform 
proves unimodality of the sequence of hook-multiplicities (Proposition~\ref{prop:s=1_unimodal}). This, in turn, implies the desired non-negativity of the quotient.  

In Section~\ref{sec:Witt} we lift the single-cycle result to the case of cycle type $(r^s)$ with $s>1$. 
In Subsection~\ref{sec:e} we provide an explicit expression for the coefficients of $(1+x)M_r(x,y)$, see Theorem~\ref{prop:e_i_formuli}.
This expression is used to obtain a product formula 
for the bivariate polynomial $1+(1+x)M_r(x,y)$
in Subsection~\ref{sec:prod}, 
and to deduce Theorem~\ref{t:main_hook-alternating2} in Subsection~\ref{sec:nn}. 

Additional results are presented in Section~\ref{sec:final}. 
In Subsection~\ref{remarks_on_s=1} it is shown that 
Lemma~\ref{lem:char_eval_cycles} and Theorem~\ref{prop:e_i_formuli} 
imply well-known combinatorial identities. 
In Subsection~\ref{sec:Cellini} it is shown that the natural approach does not provide a cyclic descent extension for conjugacy classes in $S_n$.
Palindromicity of the hook-multiplicity generating function $M_{(r^s)}(x)$ is studied in Subsection~\ref{sec:palindrom}. 

Section~\ref{sec:open} concludes the paper with final remarks and open problems.


\section{Preliminaries}\label{sec:background}
	

The role of quasisymmetric functions in the study of the distribution of descent 
sets is discussed in Subsection~\ref{sec:prel1}. 
The results presented here are 	used in Section~\ref{sec:hooks} 
to establish 
the reduction of existence of cyclic descent extension on conjugacy classes to the study of hook-multiplicities in higher Lie characters.
In Subsection~\ref{sec:prel2} we  
present cyclic analogues of a few classical results. These analogues are used for  
certain enumerative applications; the reader may skip this subsection.
	

\subsection{Quasisymmetric functions and descents}\label{sec:prel1} 

A symmetric function is called {\em Schur-positive} if all the
coefficients in its expansion in the basis of Schur functions are non-negative. Recall the notation 
$s_{\lambda/\mu}$ for the Schur function indexed by a skew shape $\lambda/\mu$. 
	
\begin{defn}\label{def:fundamental_QSF}
	For each subset $D \subseteq [n-1]$ define the {\em
		fundamental quasisymmetric function}
	\[
	   \f_{n,D}({\bf x}) := \sum_{\substack{i_1\le i_2 \le \ldots \le i_n \\ {i_j < i_{j+1} \text{ if } j \in D}}} 
	   x_{i_1} x_{i_2} \cdots x_{i_n}.
	\]
\end{defn}
 		
Denote the set of standard Young tableaux of skew shape $\lambda/\mu$ by $\SYT(\lambda/\mu)$.
There is an established notion of descent set for $\SYT(\lambda/\mu)$
\begin{equation}\label{def:Des_SYT}
    \Des(T) := 
    \{1 \le i \le n-1 \,:\, i+1 \text{ appears in a lower row of }T\text{ than }i\}.
\end{equation}

\begin{theorem}{\rm (Gessel, see~\cite[Theorem 7.19.7]{EC2})}\label{G1} 
    For every skew shape $\lambda/\mu$,
    \[
        \sum\limits_{T\in \SYT(\lambda/\mu)}\f_{n,\Des(T)}=s_{\lambda/\mu}.
    \]
\end{theorem}	
		
Given any subset $\A\subseteq S_n$, define the quasisymmetric function
\[
    \Q(\A) := \sum_{a \in \A} \f_{n,\Des(a)}.
\]

Finding subsets of permutations $\A\subseteq S_n$, for which $\Q(\A)$ is symmetric (Schur-positive), is a  long-standing problem, see~\cite{GR}. 

We write $\lambda\vdash n$ to denote that $\lambda$ is a partition of a positive integer $n$.  
For $D\subseteq [n-1]$ let ${\bf x}^D:=\prod\limits_{i\in D}x_i$.

\begin{lemma}\label{lem:Schur-gf}
    For every subset $\A\subseteq S_n$ and a family $\{c_\lambda\}_{\lambda \vdash n}$
    of coefficients, the equality
    \begin{equation}\label{eq:l1}
        \Q(\A) =\sum_{\lambda\vdash n} c_{\lambda} s_{\lambda}
    \end{equation}
    is equivalent to the equality
    \begin{equation}\label{eq:l2}
        \sum\limits_{a\in \A} {\bf x}^{\Des(a)} = \sum_{\lambda\vdash n} c_{\lambda} 
        \sum_{T\in \SYT(\lambda)}{\bf x}^{\Des(T)}.
    \end{equation}
\end{lemma}

\begin{proof}
    By Theorem~\ref{G1}, 
    Equation~\eqref{eq:l1} is equivalent to
    \[
        \sum\limits_{a\in \A} \f_{n,\Des(a)}= \sum\limits_{\lambda\vdash n} c_{\lambda} \sum\limits_{T\in \SYT(\lambda)}\f_{n,\Des(T)}.
    \]
    Next recall from~\cite[Ch.~7]{EC2} that the fundamental quasisymmetric functions in $x_1,\dots,x_n$
    form a basis of the vector space ${\rm{QSym}}_n$ of quasisymmetric functions in $n$ variables.
    Finally, apply the vector space isomorphism from ${\rm{QSym}}_n$ 
    to the space of 
    square-free polynomials in $x_1,\ldots,x_n$, which maps 
    $\f_{n,D}$ to ${\bf x}^D$. 
\end{proof}

\begin{corollary}\label{cor:Schur-gf}
    For every finite family $S$ of skew shapes of size $n$  and every subset $\A\subseteq S_n$ 
    $\Q(\A)= \sum\limits_{\lambda/\mu\in S} c_{\lambda/\mu} s_{\lambda/\mu}$ 
    if and only if 
    \[
        \sum\limits_{a\in \A} {\bf x}^{\Des(a)}=\sum\limits_{\lambda/\mu\in S} c_{\lambda/\mu} 
        \sum\limits_{T\in \SYT(\lambda/\mu)}{\bf x}^{\Des(T)}. 
    \]
\end{corollary}

Corollary~\ref{cor:Schur-gf} will be combined with Theorem~\ref{conj1} to provide criteria for the existence of cyclic descent extensions for Schur-positive sets; see the proof of Lemma~\ref{lem:2} and Remark~\ref{rem:distinct} below.  

Let $\lambda \vdash n$ be a partition of $n$ and let $\psi^\lambda$ be {\em the higher Lie character} indexed by $\lambda$ (see Definition~\ref{def:higher}).
The following result was proved by Gessel and Reutenauer. 

\begin{theorem}[{\cite[Proof of Theorem 2.1]{GR}}]\label{thm:GR22}
	For every partition $\lambda$ of $n\ge 1$, 
	\[
		\Q(\C_\lambda) = \ch(\psi^\lambda),
	\]
	where $\ch$ is the Frobenius characteristic map.
    Equivalently,		
	\[
		\Q(\C_\lambda)=
		\sum\limits_{\mu \vdash n} 
		\langle \psi^\lambda, \chi^\mu  \rangle s_\mu.
	\]
	In particular, $\Q(\C_\lambda)$ is Schur-positive.
\end{theorem}


\subsection{Cyclic analogues}
\label{sec:prel2}
		
Recall the complete homogeneous symmetric functions
\[
    h_n := 
    \sum_{i_1\le \cdots\le i_n} x_{i_1} \cdots x_{i_n}
    \qquad (n \ge 1).
\]
For a composition $\alpha =(\alpha_1,\ldots,\alpha_t)$, define
\[
    h_\alpha :=
    h_{\alpha_1}h_{\alpha_2} \cdots h_{\alpha_t}. 
\]
For any subset $J=\{j_1  < \ldots < j_t\} \subseteq [n-1]$, define
\[
    \comp{n}{J} := 
    (j_1,j_2-j_1,j_3-j_2,\ldots,j_t-j_{t-1},n-j_t).
\]
This is a composition of $n$, with a corresponding {\em connected ribbon} having the entries of $\alpha(J,n)$ as row lengths, from bottom to top. 
The associated {\em ribbon Schur function} is
\[
    s_{\alpha(J,n)} :=
    \sum_{I \subseteq J} (-1)^{\#(J \setminus I)} h_{\comp{n}{I}}. 
\]

\begin{theorem} [{Gessel, an immediate consequence of~\cite[Theorem 3]{Gessel}}]\label{thm:fiber-Des}
    Let  $\A$ be a finite set, equipped with a descent map
    $\Des: \A \longrightarrow 2^{[n-1]}$. 
    If 
    \[
        \Q(\A) :=
        \sum_{a\in \A} \f_{n,\Des(a)}
    \]
    is symmetric then
    \[
        |\{a \in \A \,:\, \Des(a)=J\}| = \langle \Q(A), s_{\alpha(J,n)} \rangle
        \qquad (\forall J\subseteq [n-1]).
    \]
\end{theorem}

For a subset $\varnothing \ne J = \{ j_1 < j_2 < \ldots < j_t \}
\subseteq [n]$
define the corresponding {\em cyclic composition} of $n$  as
\[
    \ccomp{n}{J} := 
    (j_2-j_1, \ldots, j_t - j_{t-1}, j_1 + n - j_t),
\]
with $\ccomp{n}{J} := (n)$ when $J = \{j_1\}$;
note that $\ccomp{n}{\varnothing}$ is not defined.
The corresponding {\em affine (cyclic) ribbon Schur function} was defined in~\cite{ARR} as
\[
    \cs_{\ccomp{n}{J}} := 
    \sum_{\varnothing \neq I \subseteq J}
    (-1)^{\#(J \setminus I)} h_{\ccomp{n}{I}}.
\]

\begin{theorem}[{\cite[Cor. 4.13]{AGRR}}]\label{thm:fiber-cDes}
    Let $\A$ be a finite set, equipped with a descent map
    $\Des: \A \longrightarrow 2^{[n-1]}$ 
    which has a cyclic  extension. If
    \[
        \Q(\A):=\sum\limits_{a\in \A} \f_{n,\Des(T)}
    \]
    is symmetric then the fiber sizes of (any) cyclic descent map satisfy
    \[
        |\{a \in \A \,:\, \cDes(a) =J\}|
        = \langle \Q(\A), \cs_{\ccomp{n}{J}} \rangle
        \qquad \left( \forall \varnothing \subsetneq J \subsetneq [n] \right).
    \]
\end{theorem}

\begin{proposition}[{\cite[Lemma 2.2]{ARR}}] If $A\subseteq S_n$ carries a cyclic descent extension, then 
the cyclic descent set generating function is uniquely determined.
\end{proposition}

\begin{corollary}\label{cor:cDes-dist-Schur}
    Let $\A\subseteq S_n$ be a symmetric set which carries a cyclic descent extension 
    and $S$ be a finite set of skew shapes of size $n$ which are not connected ribbons. Then 
    for every cyclic descent extension 
    the following equations are equivalent:
    \begin{equation}\label{eq:dist1}
        \Q(\A) = \sum\limits_{\lambda/\mu\in S} c_{\lambda/\mu} s_{\lambda/\mu},
    \end{equation}
    and
    \begin{equation}\label{eq:dist2}
        \sum\limits_{a\in \A} {\bf x}^{\cDes(a)} = \sum\limits_{\lambda/\mu\in S} c_{\lambda/\mu} 
        \sum\limits_{T\in \SYT(\lambda/\mu)} {\bf x}^{\cDes(T)}.
    \end{equation}
\end{corollary}

\begin{proof}
    By Theorem~\ref{thm:fiber-cDes} and Theorem~\ref{G1},
    if Equation~\eqref{eq:dist1} holds then,
    for every $\varnothing \subsetneq J \subsetneq [n]$,
    \begin{align*}
        |\{a \in \A \,:\, \cDes(a) = J\}| 
        &= \langle \Q(\A), \cs_{\ccomp{n}{J}} \rangle
        = \langle \sum\limits_{\lambda/\mu\vdash n} c_{\lambda/\mu} s_{\lambda/\mu}, \cs_{\ccomp{n}{J}} \rangle \\
        &= \sum\limits_{\lambda/\mu\vdash n} c_{\lambda/\mu} \langle \Q(\SYT(\lambda/\mu)), \cs_{\ccomp{n}{J}} \rangle \\
        &= \sum\limits_{\lambda/\mu\vdash n} c_{\lambda/\mu} |\{T \in \SYT(\lambda/\mu) \,:\, \cDes(T) =J\}|.
    \end{align*}
    Thus Equation~\eqref{eq:dist2} holds.

    For the opposite direction, let $x_n=1$ in Equation~\eqref{eq:dist2} and apply Corollary~\ref{cor:Schur-gf} to deduce  Equation~\eqref{eq:dist1}.
\end{proof}

Theorem~\ref{thm:GR22} and Theorem~\ref{thm:fiber-Des} imply the following. 

\begin{theorem}[{\cite[Theorem 2.1]{GR}}]\label{thm:GR}
    For every 
    conjugacy class $\C_\lambda$ of cycle type $\lambda \vdash n$ the
    descent set map $\Des$
    has fiber sizes given by
    \[
        |\{\pi \in \C_\lambda :\ 
        \Des(\pi)=J\}| = \langle \Q(\C_\lambda),  s_{\alpha(J,n)}\rangle
        \qquad (\forall J\subseteq [n-1]).
    \]
\end{theorem}

The following cyclic analogue of Theorem~\ref{thm:GR} results from Theorem~\ref{thm:GR22} and Theorem~\ref{thm:fiber-cDes}.
	
\begin{theorem}\label{thm:main22}
	For every conjugacy class $\C_\lambda$, which carries a cyclic descent set extension,
	all cyclic extensions of the descent set map $\Des$ 
	have fiber sizes given by
	\[
		|\{\pi\in \C_\lambda :\ \cDes(\pi)=J\}|=\langle \Q(\C_\lambda), \cs_{\ccomp{n}{J}} \rangle \qquad (\forall \varnothing \subsetneq J\subsetneq [n]).
	\]
\end{theorem}

		
\section{The role of hooks}\label{sec:hooks}

\subsection{Hooks and near-hooks}\label{sec:near_hooks}

It turns out that hooks and near-hooks play a crucial role in the study of cyclic descent extensions.

A {\em hook} is a partition with at most one part larger than $1$. Explicitly, it has the form $(n-k,1^k)$ for some $0 \le k \le n-1$. 
A {\em near-hook} of size $n$ is a hook of size $n+1$ with its (northwestern) corner cell removed; see~\cite{AER} for a somewhat more inclusive definition of this notion. 
Equivalently, recall the {\em direct sum} operation on shapes (partitions), denoted $\lambda \oplus \mu$, 
yielding a skew shape having the diagram of $\lambda$ strictly southwest of the diagram of $\mu$, with no rows or columns in common.  
A near-hook of size $n$ is the direct sum of a one-column partition $(1^k)$ and a one-row partition $(n-k)$, for some $0 \le k \le n$. 
